\documentclass[10pt,a4paper]{amsart}
\usepackage[margin=19mm]{geometry}
\usepackage{amsmath,amssymb,mathtools}
\usepackage[scr=rsfs,cal=euler]{mathalfa}
\usepackage{xfrac}
\usepackage[unicode,hidelinks,bookmarks=false]{hyperref}
\newcommand{\ie}{i.e.\ }
\newcommand{\cf}{cf.\ }
\newcommand{\resp}{respectively\ }
\newcommand{\Subsection}{\subsection}
\providecommand{\nobreakdash}{\nobreak\hbox{-}}
\setlength{\emergencystretch}{2em}
\multlinegap0pt
\usepackage{bm}
\usepackage{bbm}
\usepackage{dsfont}
\usepackage{xspace}
\usepackage{cancel}
\usepackage{mathtools}
\usepackage{enumitem}
\usepackage{amsthm}
\newtheorem{definition}{Definition}
\newtheorem{theorem}{Theorem}
\title[Exact Traces for Vector-Valued Morrey Spaces]{Exact Traces for Vector-Valued Morrey Spaces (Theorem 3)}
\author{Rishad Shahmurov, Veli Shahmurov}
\date{Source: arXiv:2609.29148v1 (CC BY 4.0)}
\begin{document}
\maketitle

\noindent\emph{Source and licence.} Exact Traces for Vector-Valued Morrey Spaces. Version arXiv:2609.29148v1 (CC BY 4.0), distributed under the Creative Commons Attribution 4.0 International licence, as recorded in the arXiv licence metadata for this exact version and shown on the abstract page https://arxiv.org/abs/2609.29148v1. Full rights notice: licenses/arxiv-2609.29148-CC-BY-4.0.txt.\par\smallskip

\noindent\textit{Editorial scope.} Counted unit: the Theorem printed as Theorem 3 of the article (source0.tex lines 515--545, source label \texttt{thm:exact-morrey-trace-space}), delivered together with the article's own complete original proof of it (source0.tex lines 547--646, one \texttt{proof} environment of 100 lines). No mathematics is written, repaired or re-derived: statement and proof are the article's own text line for line. Required context retained uncounted: def:morrey-evolution-class (lines 272--294); thm:morrey-trace-evolution (lines 298--345). Every definition, display and label of those lines is retained unchanged. Omitted from this package and not redistributed: 1--271, 295--297, 346--514, 546--546, 647--1179. Nothing inside a retained excerpt is omitted or altered: every retained body is delivered exactly as the article's lines are written, so no omission receipt of any kind accompanies this package. Citations: the retained excerpts carry no citation command at all, so no bibliography entry of the article is transcribed and no reference is redistributed. Reference adaptation: none was needed and none was made; every reference inside the delivered unit resolves inside it, so no unresolved reference or citation is produced. Printed slips kept literal: no printed word, symbol, hypothesis or number of the counted statement or of its proof was altered. Layout and class: the article's own class is \texttt{svjour3} with packages including fix-cm, fontenc, inputenc, amsmath, amssymb, mathtools, mathrsfs, bm, enumitem, microtype, hyperref, cleveref; the wrapper is the lane's pinned \texttt{amsart} editorial wrapper at 10pt on A4, which restates the theorem-like environments the retained text needs and adds the article's own macro definitions that the retained text needs (none), with extra wrapper packages (bm, bbm, dsfont, xspace, cancel, mathtools, enumitem). The delivered numbering is the unit's own. Assets: no retained span contains a figure environment, a graphics-inclusion command, a PDF-page inclusion, a table or a TikZ picture; no image file is copied into this package, the package declares no asset dependency and no image file is redistributed with it. Licence: version arXiv:2609.29148v1 of this article is distributed under the Creative Commons Attribution 4.0 International licence, as recorded in the arXiv licence metadata for this exact version and shown on the abstract page https://arxiv.org/abs/2609.29148v1. A full rights notice is delivered at licenses/arxiv-2609.29148-CC-BY-4.0.txt. Characters: every retained excerpt is pure ASCII, so the delivered engine, XeLaTeX with its default Latin Modern text font, reports no missing character.
	\begin{definition}\label{def:morrey-evolution-class}
		Let $(X_0,X_1)$ be a Banach couple with continuous embedding
		\[
		X_1\hookrightarrow X_0.
		\]
		Fix
		\[
		1<p<\infty,\qquad 0\le \lambda<1,\qquad T>0,
		\qquad
		\theta:=1-\frac{1-\lambda}{p}\in(0,1).
		\]
		We define
		\[
		\mathbb E^{p,\lambda}(0,T;X_0,X_1)
		:=
		\Bigl\{
		u:\;
		u\in \mathcal M^{p,\lambda}(0,T;X_1),\;
		\partial_t u\in \mathcal M^{p,\lambda}(0,T;X_0)
		\Bigr\},
		\]
		where $\partial_t u$ is understood in the sense of $X_0$-valued distributions.
	\end{definition}

\begin{theorem}\label{thm:morrey-trace-evolution}
		Let $(X_0,X_1)$, $p$, $\lambda$, $T$, and $\theta$ be as in Definition~\ref{def:morrey-evolution-class}. Let
		\[
		u\in \mathbb E^{p,\lambda}(0,T;X_0,X_1).
		\]
		Then $u$ admits an $X_0$-valued absolutely continuous representative on $[0,T]$, still denoted by $u$, such that:
		
		\begin{enumerate}[label=\textup{(\roman*)}]
			\item for every $0\le s<t\le T$,
			\[
			u(t)-u(s)=\int_s^t \partial_t u(\tau)\,d\tau
			\qquad\text{in }X_0,
			\]
			and
			\[
			\|u(t)-u(s)\|_{X_0}
			\le
			|t-s|^\theta
			\|\partial_t u\|_{\mathcal M^{p,\lambda}(0,T;X_0)};
			\]
			
			\item for every $t\in[0,T]$,
			\[
			u(t)\in (X_0,X_1)_{\theta,\infty},
			\]
			and
			\[
			\sup_{t\in[0,T]}
			\|u(t)\|_{(X_0,X_1)_{\theta,\infty}}
			\le
			C\Bigl(
			\|u\|_{\mathcal M^{p,\lambda}(0,T;X_1)}
			+
			\|\partial_t u\|_{\mathcal M^{p,\lambda}(0,T;X_0)}
			\Bigr);
			\]
			
			\item in particular,
			\[
			u\in C^{0,\theta}([0,T];X_0)\cap L^\infty\bigl(0,T;(X_0,X_1)_{\theta,\infty}\bigr),
			\]
			and the endpoint traces
			\[
			u(0),\qquad u(T)
			\]
			are well-defined elements of $(X_0,X_1)_{\theta,\infty}$.
		\end{enumerate}
	\end{theorem}

\begin{theorem}\label{thm:exact-morrey-trace-space}
	Let \((X_0,X_1)\) be a Banach couple with continuous embedding \(X_1\hookrightarrow X_0\).  Let
	\[
	1<p<\infty,\qquad 0<\lambda<1,\qquad T>0,
	\qquad
	\theta=1-\frac{1-\lambda}{p}.
	\]
	Then the trace at \(t=0\) of the Morrey evolution class is exactly
	\[
	\operatorname{Tr}_{0}\mathbb E^{p,\lambda}(0,T;X_0,X_1)
	=
	(X_0,X_1)_{\theta,\infty}.
	\]
	More precisely, the trace map
	\[
	\operatorname{Tr}_{0}:u\mapsto u(0)
	\]
	is bounded from \(\mathbb E^{p,\lambda}(0,T;X_0,X_1)\), equipped with its natural norm, into \((X_0,X_1)_{\theta,\infty}\).  Conversely, for every \(x\in (X_0,X_1)_{\theta,\infty}\) there exists
	\[
	u\in \mathbb E^{p,\lambda}(0,T;X_0,X_1)
	\]
	such that \(u(0)=x\) in \(X_0\), and
	\[
	\|u\|_{\mathcal M^{p,\lambda}(0,T;X_1)}
	+
	\|u'\|_{\mathcal M^{p,\lambda}(0,T;X_0)}
	\le
	C\|x\|_{(X_0,X_1)_{\theta,\infty}}.
	\]
	The constant \(C\) depends only on \(p,\lambda,T\), and the embedding constant of \(X_1\hookrightarrow X_0\).
	\end{theorem}

	\begin{proof}
	The boundedness of the trace map is exactly the trace estimate in Theorem~\ref{thm:morrey-trace-evolution}.  We prove the converse by an explicit dyadic extension.

	Let
	\[
	M:=\|x\|_{(X_0,X_1)_{\theta,\infty}}.
	\]
	For \(n=0,1,2,\ldots\), put
	\[
	r_n:=T2^{-n}.
	\]
	By the definition of the \((X_0,X_1)_{\theta,\infty}\)-norm, for each \(n\) we may choose a decomposition
	\[
	x=a_n+b_n,
	\qquad a_n\in X_0,
	\qquad b_n\in X_1,
	\]
	such that
	\[
	\|a_n\|_{X_0}+r_n\|b_n\|_{X_1}
	\le
	2K(r_n,x;X_0,X_1)
	\le
	2Mr_n^\theta.
	\]
	Therefore
	\[
	\|x-b_n\|_{X_0}=\|a_n\|_{X_0}\le 2Mr_n^\theta,
	\qquad
	\|b_n\|_{X_1}\le 2Mr_n^{\theta-1}.
	\]
	In particular \(b_n\to x\) in \(X_0\).

	Define \(u\) on each dyadic interval \((r_{n+1},r_n]\) by linear interpolation between \(b_{n+1}\) and \(b_n\):
	\[
	u(t):=
	\frac{t-r_{n+1}}{r_n-r_{n+1}}b_n
	+
	\frac{r_n-t}{r_n-r_{n+1}}b_{n+1},
	\qquad r_{n+1}<t\le r_n.
	\]
	Set \(u(0):=x\).  Since \(b_n\to x\) in \(X_0\), the resulting function is continuous at \(0\) as an \(X_0\)-valued function.  For \(t\in(r_{n+1},r_n]\), the preceding estimates give
	\[
	\|u(t)\|_{X_1}
	\le
	\max\{\|b_n\|_{X_1},\|b_{n+1}\|_{X_1}\}
	\le
	C M r_n^{\theta-1}.
	\]
	Since \(t\simeq r_n\) on \((r_{n+1},r_n]\), and \(\theta-1=-(1-\lambda)/p\), we have
	\[
	\|u(t)\|_{X_1}
	\le
	C M t^{-\frac{1-\lambda}{p}}.
	\]
	Similarly, on \((r_{n+1},r_n)\),
	\[
	u'(t)=\frac{b_n-b_{n+1}}{r_n-r_{n+1}}.
	\]
	Using \(b_n-x=-a_n\) and \(b_{n+1}-x=-a_{n+1}\),
	\[
	\|b_n-b_{n+1}\|_{X_0}
	\le
	\|a_n\|_{X_0}+\|a_{n+1}\|_{X_0}
	\le
	C M r_n^\theta.
	\]
	Since \(r_n-r_{n+1}=r_{n+1}\simeq r_n\), it follows that
	\[
	\|u'(t)\|_{X_0}
	\le
	C M r_n^{\theta-1}
	\le
	C M t^{-\frac{1-\lambda}{p}}.
	\]

	It remains only to observe that the scalar function
	\[
	g(t)=t^{-\frac{1-\lambda}{p}},\qquad 0<t<T,
	\]
	belongs to \(\mathcal M^{p,\lambda}(0,T)\) when \(0<\lambda<1\).  Indeed, if \(I=(a,b)\subset(0,T)\) and \(h=b-a\), then
	\[
	\int_I g(t)^p\,dt=\int_a^b t^{-(1-\lambda)}\,dt.
	\]
	If \(a\le h\), then this is bounded by \(C h^\lambda\).  If \(a>h\), then
	\[
	\int_a^b t^{-(1-\lambda)}\,dt
	\le
	h a^{-(1-\lambda)}
	\le
	h h^{-(1-\lambda)}=h^\lambda.
	\]
	Thus
	\[
	|I|^{-\lambda/p}
	\left(\int_I g(t)^p\,dt\right)^{1/p}
	\le C,
	\]
	uniformly in \(I\).  Hence \(u\in\mathcal M^{p,\lambda}(0,T;X_1)\) and \(u'\in\mathcal M^{p,\lambda}(0,T;X_0)\), with the asserted estimate.
	\end{proof}

\end{document}
